% from ega4-2-fr.tex
\oldpage[IV-2]{117}

formons pour ces deux composantes une suite $(X_i)$ ayant la propriété énoncée
dans b) et posons $U_i=X_i\cap U$ pour $1\leq i\leq n$; alors $U_i$ est une
composante irréductible de $U$ (0$_\mathrm I$, 2.1.6) et de plus
$U_i\cap U_{i-1}\neq\varnothing$ pour $1\leq i\leq n$, sinon on aurait
$X_i\cap X_{i-1}\subset Y$, donc $X_i\cap X_{i-1}\in\mathfrak F$,
contrairement à la définition des $X_i$. Cela entraîne que $U$ est connexe.

Montrons maintenant que a) entraîne b). Désignons par $Y$ la réunion de la
famille $(X_\alpha\cap X_\beta)$, où $(X_\alpha,X_\beta)$ parcourt l'ensemble
des couples de composantes irréductibles distinctes de $X$ telles que
$X_\alpha\cap X_\beta\in\mathfrak F$. Pour tout point $y\in Y$, il y a un
voisinage ouvert $V$ de $y$ dans $X$ ne rencontrant qu'un nombre fini de
composantes irréductibles de $X$; cela montre d'une part que $Y$ est fermé et
d'autre part que $V\cap Y$ est intersection de $V$ et d'un ensemble de
$\mathfrak F$; en vertu de l'hypothèse faite sur $\mathfrak F$, on a
$Y\in\mathfrak F$, donc $U=X-Y$ est connexe; en outre $Y$ est rare dans $X$.
Soient alors $X'$, $X''$ deux composantes irréductibles distinctes de $X$,
$U'$, $U''$ leurs traces respectives sur $U$; ce sont des composantes
irréductibles distinctes de $U$ (0$_\mathrm I$, 2.1.6). Or la réunion des
composantes irréductibles $W$ de $U$ telles qu'il existe une suite
$(U_i)_{0\leq i\leq n}$ de composantes irréductibles de $U$ pour laquelle
$U_0=U'$, $U_{i-1}\neq U_i$ et $U_{i-1}\cap U_i\neq\varnothing$ pour
$1\leq i\leq n$ et $U_n=W$ est un ensemble ouvert et fermé dans $U$, puisque
$U$ est localement noethérien et par suite ses composantes irréductibles
forment une famille localement finie d'ensembles fermés. Il y a donc une telle
suite $(U_i)$ pour laquelle $U_n=U''$; soit $X_i$ ($0\leq i\leq n$) la
composante irréductible de $X$ telle que $X_i\cap U=U_i$
(0$_\mathrm I$, 2.1.6); comme $U_{i-1}\neq U_i$, on a
$X_{i-1}\neq X_i$ pour $1\leq i\leq n$; si l'on avait
$X_{i-1}\cap X_i\in\mathfrak F$ pour un $i$ tel que $1\leq i\leq n$, on en
déduirait $X_{i-1}\cap X_i\subset Y$ par définition de $Y$, d'où
$U_{i-1}\cap U_i=\varnothing$, contrairement à l'hypothèse. Ceci achève la
démonstration du lemme.

On notera que l'hypothèse faite sur $X$ dans (5.10.8) est vérifiée lorsque $X$
est un préschéma de Cohen-Macaulay et $d\geq2$.

\begin{corollary}[5.10.9]
\label{IV.2.s5.10.9-fr}
Si un anneau local noethérien $A$ vérifie $(S_2)$ et est caténaire, il est
équidimensionnel.
\end{corollary}

L'hypothèse de (5.10.8) est alors vérifiée par $X=\operatorname{Spec}(A)$,
avec $d=2$. Pour montrer que toutes les composantes irréductibles de $X$ ont
même dimension, il suffit alors, en vertu de (5.10.8), de montrer que deux
telles composantes $X'$, $X''$ ont même dimension lorsque l'on suppose en
outre que $\operatorname{codim}(X'\cap X'',X)=1$. Il y a alors une composante
irréductible $Z$ de $X'\cap X''$ telle que
$\operatorname{codim}(Z,X)=1$, donc $\operatorname{codim}(Z,X')=1$, puisque
$\operatorname{codim}(Z,X)\geq\operatorname{codim}(Z,X')\geq1$; de même
$\operatorname{codim}(Z,X'')=1$, et comme $X$ est caténaire, cela entraîne
$\dim(X')=\dim(X'')$.

\begin{proposition}[5.10.10]
\label{IV.2.s5.10.10-fr}
Soient $X$ un préschéma localement noethérien, $Z$ une partie de $X$ stable
par spécialisation, $\mathcal F$ un $\mathcal O_X$-Module cohérent et supposons
que le $\mathcal O_X$-Module $\mathcal H^0_{X/Z}(\mathcal F)$ soit cohérent.
Alors:
\begin{enumerate}
\item[(i)] On a
$\operatorname{prof}_Z(\mathcal H^0_{X/Z}(\mathcal F))\geq2$.
\item[(ii)] Pour tout point
$x\in\operatorname{Ass}(\mathcal F)\cap(X-Z)$, on a
\[
\operatorname{codim}(Z\cap\overline{\{x\}},\overline{\{x\}})\geq2.
\]
\item[(iii)] L'ensemble $U$ des $x\in X$ tels que
$\operatorname{prof}_{Z_x}(\widetilde{\mathcal F}_x)\geq2$ (notations de
(5.9.6)) est ouvert dans $X$; on a $X-U\subset Z$, et $U$ est le plus grand
ouvert de $X$ tel que $\mathcal F|U$ soit $(Z\cap U)$-clos.
\end{enumerate}
\end{proposition}

